\chapter{Реестры и доставка образов}

Собранный образ нужно доставить туда, где он будет работать:
на другой компьютер, в тестовый стенд, в production. Для этого служат
реестры~--- хранилища образов с версионированием и контролем доступа.
В главе~--- публикация, теги, мультиархитектурные сборки, собственный
реестр и автоматизация в CI.

\section{Что такое реестр}

\begin{defbox}
\textbf{Реестр} (\emph{registry})~--- сервис, хранящий образы и
предоставляющий API для их загрузки и выгрузки. Реестр хранит слои
по digest, поэтому одинаковые слои не дублируются.
\end{defbox}

Популярные реестры:

\begin{center}
\small
\begin{tabularx}{\textwidth}{@{}lX@{}}
\toprule
\textbf{Реестр} & \textbf{Особенности} \\
\midrule
Docker Hub & публичный по умолчанию, огромная библиотека официальных
образов, лимиты на анонимные загрузки \\
GitHub Container Registry (GHCR) & удобен для проектов на GitHub,
токен \cmd{GITHUB\_TOKEN} \\
GitLab Container Registry & встроен в GitLab, права наследуются
из проекта \\
Amazon ECR, Google Artifact Registry, Azure ACR & облачные реестры
с интеграцией IAM \\
Harbor & self-hosted, сканирование, политики, репликация \\
registry:2 & минимальный эталонный реестр для собственных нужд \\
\bottomrule
\end{tabularx}
\end{center}

С точки зрения Docker все реестры выглядят одинаково: имя образа
начинается с домена реестра, а команды \cmd{pull} и \cmd{push}
работают по одному протоколу.

\section{Аутентификация}

Публичные образы можно скачивать без входа. Для публикации и для
приватных образов нужна аутентификация.

\begin{code}[language=bash]
$ docker login
Login with your Docker ID to push and pull images from Docker Hub.
Username: acme
Password:
Login Succeeded

$ docker login ghcr.io -u acme
$ docker logout
\end{code}

\begin{tipbox}
В CI используйте токены и флаг \flag{--password-stdin}:
\cmd{printf '\%s' "\$TOKEN" | docker login ghcr.io -u USER --password-stdin}.
Так пароль не попадает в историю команд и списки процессов.
\end{tipbox}

Учётные данные сохраняются в \file{\textasciitilde/.docker/config.json},
обычно через credential helper (например, \cmd{pass} или
\cmd{secretservice}). Если helper не настроен, пароль может лежать
в файле в base64~--- это не шифрование.

\begin{warnbox}
Не публикуйте \file{\textasciitilde/.docker/config.json} и не передавайте
его в контейнеры. Для CI создавайте отдельные токены с минимальными
правами и отзывайте их после использования.
\end{warnbox}

\section{Имена, теги и стратегии версионирования}

При публикации имя образа должно содержать домен реестра и пространство
имён:

\begin{code}[language=bash]
$ docker tag demo-api:1.0 ghcr.io/acme/demo-api:1.0
$ docker push ghcr.io/acme/demo-api:1.0
\end{code}

Один образ можно отметить несколькими тегами. Разумная стратегия:

\begin{itemize}
  \item \textbf{неизменяемая версия}~--- \cmd{1.4.2} (semver) или
        \cmd{sha-3f2a1c7} (короткий хеш коммита);
  \item \textbf{скользящая минорная}~--- \cmd{1.4} (обновляется
        с патчами);
  \item \textbf{среда}~--- \cmd{staging}, \cmd{production} (указывает
        на конкретную сборку, но подвижна);
  \item \cmd{latest}~--- только для удобства разработки, не для
        production.
\end{itemize}

\begin{code}[language=bash]
$ SHA=$(git rev-parse --short HEAD)
$ docker build -t ghcr.io/acme/demo-api:$SHA -t ghcr.io/acme/demo-api:1.4.2 .
$ docker push ghcr.io/acme/demo-api:$SHA
$ docker push ghcr.io/acme/demo-api:1.4.2
$ docker push --all-tags ghcr.io/acme/demo-api
\end{code}

\begin{notebox}
Многие реестры поддерживают неизменяемые теги (tag immutability):
попытка перезаписать тег отклоняется. Если такой возможности нет,
фиксируйте digest в манифестах развёртывания~--- это единственная
по-настоящему неизменяемая ссылка.
\end{notebox}

\section{Мультиархитектурные образы}

Ноутбуки на ARM и серверы на x86 требуют разных сборок. BuildKit
позволяет собрать и опубликовать один тег, указывающий на несколько
платформ.

\begin{code}[language=bash]
$ docker buildx create --name multi --use --driver docker-container
$ docker buildx build \
    --platform linux/amd64,linux/arm64 \
    --tag ghcr.io/acme/demo-api:1.4.2 \
    --push .
$ docker buildx imagetools inspect ghcr.io/acme/demo-api:1.4.2
Manifests:
  Platform: linux/amd64
  Platform: linux/arm64
\end{code}

\begin{warnbox}
Мультиплатформенный результат нельзя загрузить в локальный Docker
флагом \flag{--load}: локальное хранилище поддерживает только одну
платформу. Публикуйте через \flag{--push} или собирайте под свою
платформу отдельно.
\end{warnbox}

\begin{tipbox}
Для сборки под чужую архитектуру на своей машине Docker использует
эмуляцию QEMU. Она подходит для проверки, но сборка может быть
в разы медленнее, а некоторые программы ведут себя иначе. В CI
используйте отдельные runner'ы под каждую платформу.
\end{tipbox}

\section{Собственный реестр}

Для локальной сети или закрытого контура можно поднять реестр
за пару минут.

\begin{code}[language=bash]
$ docker run -d --name registry -p 5000:5000 \
    -v registry-data:/var/lib/registry \
    --restart unless-stopped \
    registry:2

$ docker tag demo-api:1.0 localhost:5000/demo-api:1.0
$ docker push localhost:5000/demo-api:1.0
$ curl -s http://localhost:5000/v2/_catalog
{"repositories":["demo-api"]}
\end{code}

Для доступа с других машин по HTTP (без TLS) реестр нужно объявить
небезопасным в \file{/etc/docker/daemon.json} на каждой машине:

\begin{code}[language=bash]
{
  "insecure-registries": ["registry.local:5000"]
}
\end{code}

\begin{warnbox}
Небезопасный реестр допустим только в изолированной тестовой сети.
В production используйте TLS и аутентификацию. Готовые решения~---
Harbor, GitLab Registry, облачные реестры~--- уже содержат HTTPS,
разграничение прав и сканирование уязвимостей.
\end{warnbox}

Базовая аутентификация через \cmd{htpasswd}:

\begin{code}[language=bash]
$ mkdir -p auth
$ docker run --rm --entrypoint htpasswd httpd:2 \
    -Bbn acme secret > auth/htpasswd

$ docker run -d --name registry -p 5000:5000 \
    -v "$(pwd)/auth:/auth" \
    -e REGISTRY_AUTH=htpasswd \
    -e REGISTRY_AUTH_HTPASSWD_REALM="Registry Realm" \
    -e REGISTRY_AUTH_HTPASSWD_PATH=/auth/htpasswd \
    registry:2
\end{code}

\section{Публикация из CI}

Типичный конвейер: собрать образ, прогнать тесты, опубликовать тег,
зафиксировать digest. Пример для GitHub Actions:

\begin{code}[language=composeyaml]
name: build-and-push
on:
  push:
    tags: ["v*"]

jobs:
  image:
    runs-on: ubuntu-latest
    permissions:
      contents: read
      packages: write
    steps:
      - uses: actions/checkout@v4
      - uses: docker/setup-buildx-action@v3
      - uses: docker/login-action@v3
        with:
          registry: ghcr.io
          username: ${{ github.actor }}
          password: ${{ secrets.GITHUB_TOKEN }}
      - uses: docker/metadata-action@v5
        id: meta
        with:
          images: ghcr.io/${{ github.repository }}
          tags: |
            type=ref,event=tag
            type=sha
      - uses: docker/build-push-action@v6
        with:
          context: .
          push: true
          tags: ${{ steps.meta.outputs.tags }}
          labels: ${{ steps.meta.outputs.labels }}
          cache-from: type=gha
          cache-to: type=gha,mode=max
\end{code}

\begin{itemize}
  \item \cmd{setup-buildx} включает BuildKit и кэш;
  \item \cmd{login-action} безопасно передаёт токен;
  \item \cmd{metadata-action} формирует теги из git-тега и хеша коммита;
  \item \cmd{build-push-action} собирает и публикует образ,
        используя кэш GitHub Actions.
\end{itemize}

\begin{tipbox}
Собирайте образ ровно один раз и продвигайте его между средами
по digest, а не пересобирайте на каждом этапе. Иначе тесты
и production могут получить разные артефакты из одного коммита.
\end{tipbox}

\section{Доверие к образам}

Публикация~--- не конец пути. Нужно понимать, что именно вы запускаете.

\begin{itemize}
  \item \textbf{Сканирование уязвимостей.} \cmd{docker scout cves IMAGE},
        \cmd{trivy image IMAGE}, встроенные сканеры Harbor и облачных
        реестров.
  \item \textbf{SBOM.} Список компонентов образа:
        \cmd{docker buildx build --sbom=true --push ...} или
        \cmd{syft IMAGE}.
  \item \textbf{Provenance.} Сведения о том, как и из чего собран образ:
        \cmd{--provenance=true}.
  \item \textbf{Подпись.} \cmd{cosign sign IMAGE}, проверка
        \cmd{cosign verify IMAGE}; политика допускает только подписанные
        образы.
  \item \textbf{Pinning.} Фиксация базовых образов по digest:
        \cmd{FROM node:22-alpine@sha256:...}.
\end{itemize}

\begin{warnbox}
Тег \cmd{nginx:1.27-alpine} может указывать на разные digest в разные
дни: сопровождающие пересобирают образ с обновлениями безопасности.
Это нормально, но означает, что «та же версия» не гарантирует тот же
байт-в-байт образ. Для строгой воспроизводимости фиксируйте digest
и обновляйте его осознанно.
\end{warnbox}

\section{Лимиты и кэширование}

Docker Hub ограничивает число скачиваний с одного IP для анонимных
пользователей и отдельно~--- для авторизованных. Лимиты меняются,
поэтому проверяйте актуальные условия. Практические меры:

\begin{itemize}
  \item авторизуйтесь в CI (\cmd{docker login})~--- лимит выше;
  \item используйте pull-through cache или зеркало реестра;
  \item в облаках включайте кэширование образов на узлах;
  \item не скачивайте образ на каждом шаге конвейера~--- один раз
        на сборку достаточно.
\end{itemize}

В Compose можно управлять политикой скачивания:

\begin{code}[language=composeyaml]
services:
  api:
    image: ghcr.io/acme/demo-api:1.4.2
    pull_policy: always
\end{code}

\begin{summarybox}
\begin{itemize}
  \item Реестр хранит образы и слои; имя образа включает домен реестра.
  \item Публикация требует аутентификации; в CI используйте токены
        и \flag{--password-stdin}.
  \item Теги подвижны, digest неизменяем; для production фиксируйте
        версию или digest.
  \item Мультиархитектурные образы собираются через
        \cmd{docker buildx} и публикуются флагом \flag{--push}.
  \item Доверие к образу обеспечивают сканирование, SBOM, provenance
        и подписи.
\end{itemize}
\end{summarybox}

\section*{Вопросы для самопроверки}
\begin{selfcheck}
  \item Зачем в имени образа указывать домен реестра?
  \item Почему \cmd{latest} не подходит для production?
  \item Чем тег отличается от digest?
  \item Почему мультиплатформенный образ нельзя загрузить флагом
        \flag{--load}?
  \item Какие меры снижают риск запуска скомпрометированного образа?
  \item Как передать токен в \cmd{docker login} без пароля в аргументах?
\end{selfcheck}
